\chapter{Bronchoscopy and Bronchoalveolar Lavage (BAL)}

\section*{What Is This Test?}
Bronchoscopy uses a flexible or rigid camera to examine the airways. Bronchoalveolar lavage (BAL) adds sterile fluid through the bronchoscope and retrieves it for cell, microbiologic, or other analysis from a selected lung region.

\section*{Alternative Names}
\begin{itemize}
  \item Flexible Bronchoscopy
  \item Fiberoptic Bronchoscopy
  \item BAL Fluid Analysis
  \item Bronchial Washing
\end{itemize}

\section*{Principle}
The bronchoscope provides direct visualization and permits washing, brushing, biopsy, foreign-body removal, and secretion sampling. BAL fluid may be examined with stains, cultures, PCR, cytology, differential cell counts, or specialized tests.

\section*{Why This Test Is Ordered}
It may investigate persistent cough, hemoptysis, abnormal imaging, airway obstruction, suspected lung cancer, difficult-to-diagnose infection, diffuse lung disease, aspiration, or a foreign body. It can also help treat airway obstruction or remove secretions.

\section*{Primary vs Secondary Test}
Bronchoscopy is an invasive secondary test selected after history, examination, imaging, and less-invasive tests. BAL results are most useful when the sampled lung region and suspected disease are clearly defined.

\section*{How This Test Is Performed}
After fasting and sedation or anesthesia, the scope passes through the nose or mouth into the airways. For BAL, sterile saline is instilled and suctioned into a collection trap. Biopsy or brushing may be performed when indicated.

\section*{What Preparation Is Needed}
Follow fasting and medication instructions. Tell the team about anticoagulants, bleeding disorders, sleep apnea, oxygen use, allergies, and pregnancy possibility. Arrange an escort after sedation.

\section*{Contraindications and Precautions}
Temporary cough, sore throat, hoarseness, fever, low oxygen, bleeding, infection, bronchospasm, and rarely pneumothorax or serious cardiopulmonary complications can occur. Biopsy increases bleeding and pneumothorax risk.

\section*{Understanding Results}
Findings may show inflammation, infection, malignant cells, airway narrowing, bleeding, or a diffuse inflammatory cell pattern. A negative BAL does not exclude disease because sampling is limited and some conditions are patchy or require tissue biopsy.

\section*{Follow-up Tests}
Follow-up may include bacterial, fungal, mycobacterial, or viral studies; cytology and biopsy review; CT or PET imaging; pulmonary-function tests; repeat bronchoscopy; or surgical lung biopsy in selected cases.

\section*{References}
\begin{enumerate}
  \item MedlinePlus. Bronchoscopy and Bronchoalveolar Lavage (BAL).
  \item American Thoracic Society. Bronchoscopy and BAL resources.
  \item American College of Chest Physicians. Diagnostic bronchoscopy guidance.
\end{enumerate}

\clearpage
\section*{Understanding Results and Follow-up}
The laboratory report should identify the specimen, method, units, reference interval or decision limit, and whether the result is positive, negative, abnormal, or inconclusive. Reference intervals vary with age, sex, pregnancy, specimen type, assay, and laboratory; a result outside a range is not automatically a diagnosis, and a result within a range does not always exclude disease. Discuss unexpected results with the ordering clinician, who can decide whether repeat or confirmatory testing, treatment, or referral is appropriate. Do not change medicines or delay urgent care based on an isolated result.


\section*{Regulatory Status and Authorization Date}
This chapter describes a test category, not a specific commercial product. FDA, CDSCO, EU IVDR, and MHRA approval or registration dates therefore cannot be assigned without the assay manufacturer, product name, instrument, and intended use. For laboratory-developed tests, an individual agency approval date may not exist. See the Preface for the interpretation of regulatory status.

\clearpage
